% Part: computability
% Chapter: recursive-functions
% Section: trees

\documentclass[../../../include/open-logic-section]{subfiles}

\begin{document}

\olfileid{cmp}{rec}{tre}
\olsection{ونې}

ځينې وخت د ونو ښودل د طبيعي عددونو په توګه ګټور وي، کټ مټ لکه لړۍ چې په عددونو ښيو او د هغو خاصيتونه او عملونه د کوډونو پر بنسټيزو بازګشتي اړيکو او تابعو ښيو۔ د دې دپاره لړۍ او د هغو کوډونه کاروو۔ يوه ونه يا يوازې يوه غوټه (چې ښايي لېبل ولري) وي، يا يوه غوټه (چې ښايي لېبل ولري) وي چې له يو شمېر فرعي ونو سره تړلې وي۔ دې غوټې ته د ونې \emph{ريښه}، او هغو فرعي ونو ته چې ورسره تړلې وي د هغې \emph{سمدستي فرعي ونې} وايي۔

ونې په بازګشتي ډول د~$\tuple{k, d_1, \dots, d_k}$ په لړۍ کوډوو، چې $k$ د سمدستي فرعي ونو شمېر او~$d_1$، \dots،~$d_k$ د سمدستي فرعي ونو کوډونه دي۔ که غوټې لېبلونه ولري، له سمدستي فرعي ونو وروسته ئې شاملولے شو۔ نو هغه ونه چې يوازې يوه د~$l$ لېبل لرونکې غوټه ولري، په~$\tuple{0,l}$ کوډېږي؛ او هغه ونه چې يوه د~$l_1$ لېبل لرونکې ريښه ئې له دوو يوازېنيو غوټو سره تړلې وي، چې لېبلونه ئې~$l_2$ او~$l_3$ دي، په~$\tuple{2, \tuple{0, l_2}, \tuple{0, l_3}, l_1}$ کوډېږي۔

\begin{prop}
  \ollabel{prop:subtreeseq}
  تابع~$\fn{SubtreeSeq}(t)$، چې د داسې لړۍ کوډ راګرځوي چې غړي ئې د $t$ کوډ لرونکې ونې د ټولو فرعي ونو کوډونه دي، بنسټيزه بازګشتي ده۔
\end{prop}

\begin{proof}
  لومړے پام وکړئ چې~$\fn{ISubtrees}(t) = \fn{subseq}(t, 1, (t)_0)$ بنسټيزه بازګشتي ده او د ونې~$t$ د سمدستي فرعي ونو کوډونه راګرځوي۔ اوس مرستندويه تابع~$\fn{hSubtreeSeq}(t,n)$ داسې تعريفوو چې د ټولو هغو فرعي ونو لړۍ محاسبه کړي چې ريښې ئې د~$t$ له ريښې څخه تر ټولو زيات~$n$ څنډې لرې وي۔ د~$t$ د هغو فرعي ونو لړۍ چې له ريښې څخه تر ټولو زيات~$0$ څنډې لرې دي---يعنې پخپله د~$t$ په ريښه پيلېږي---يوازې له~$t$ څخه جوړه لړۍ ده۔ د~$n+1$ تر کچې د ټولو فرعي ونو د لړۍ د ترلاسه کولو دپاره، د~$n$ تر کچې شته فرعي ونې د هغو ټولو سمدستي فرعي ونو له لړۍ سره نښلوو چې د~$n$ تر کچې له فرعي ونو راځي۔
  \textbf{د ژباړې د نسخې سمون (OLCMP-013):} سرچينې مرستندويه تابع له ريښې څخه په کټ مټ ټاکلي واټن د فرعي ونو لړۍ بلله، خو لاندې بازګشتي ګام مخکينۍ لړۍ هم ساتي او تر ټاکلي واټن پورې ټولې کچې راټولوي۔ دلته تشريح د هماغې معادلې له مجموعي چلند سره برابره شوې ده۔
  د ټولو د لست د ترلاسه کولو دپاره پام وکړئ: که~$f(x)$ يوه بنسټيزه بازګشتي تابع وي چې د لړيو کوډونه راګرځوي، نو د غړو د مثبت شمېر دپاره تابع~$g_f(s,k)=f((s)_0) \concat \dots \concat f((s)_{k-1})$ هم بنسټيزه بازګشتي ده؛ د صفر غړو نښلون تشه لړۍ نيسو:
    \begin{align*}
      g(s, 0) & = \emptyseq\\
      g(s, k+1) & = g(s, k) \concat f((s)_k)
    \end{align*}
    مثلاً، که~$s$ د ونو يوه لړۍ وي، نو $h(s) = g_{\fn{ISubtrees}}(s, \len{s})$ د~$s$ د غړو د ټولو سمدستي فرعي ونو لړۍ راکوي۔
    \textbf{د ژباړې د نسخې سمون (OLCMP-014):} سرچينې مرستندويه نښلون د صفرم غړي په تابع پيل کاوه او بيا ئې تر اوږدوالي پورې ټول‌شموله شاخص غځاوه؛ په دې سره ئې د لړۍ له وروستي قانوني شاخص وروسته يو زيات غړے لوسته۔ دلته پاراميټر د لومړيو غړو شمېر دے: صفر تش نښلون ورکوي او هر راتلونکی ګام سملاسي ورپسې قانوني غړے زياتوي۔
    له دې څخه~$\fn{hSubtreeSeq}$ داسې تعريفولے شو:
    \begin{align*}
      \fn{hSubtreeSeq}(t, 0) & = \tuple{t} \\
      \fn{hSubtreeSeq}(t, n+1) & = \fn{hSubtreeSeq}(t, n) \concat h(\fn{hSubtreeSeq}(t, n)).
    \end{align*}
    د~$t$ په کوډ شوې ونه کښې د فرعي ونو تر ټولو لوړه کچه، يعنې د ريښې او يوې پاڼه غوټې تر منځ تر ټولو لوے واټن، د کوډ~$t$ له مخې محدوده ده۔ نو د~$t$ په کوډ شوې ونه کښې د ټولو فرعي ونو د کوډونو لړۍ~$\fn{hSubtreeSeq}(t, t)$ ورکوي۔
\end{proof}

\begin{prob}
  د~\olref[cmp][rec][tre]{prop:subtreeseq} په ثبوت کښې د $\fn{hSubtreeSeq}$ تعريف په عمومي ډول تکرارونه راولي۔ يو داسې بل تعريف ورکړئ چې ضمانت وکړي د يوې فرعي ونې کوډ په پايله لست کښې يوازې يو ځل راځي۔
\end{prob}

\end{document}
